% chapters/01_introduction.tex
% Target: 500–600 words
% ─────────────────────────────────────────────────────────────────────────────

\section{Motivation and Background}

% The semiconductor GVC as a strategic policy issue.
% Norway's dual position: upstream raw material exporter (REC Silicon, Elkem)
% + downstream importer of finished chips (sensors, quantum research).
% Gap: standard RCA/FIR metrics miss structural position in the whole network.
% SNA as the missing analytical layer.

[Write introduction here]

% ─────────────────────────────────────────────────────────────────────────────
\section{Research Questions}

The thesis is organised around one main research question and two sub-questions.

\bigskip
\noindent\textbf{Main RQ:} \textit{What does Social Network Analysis reveal about the structure of the global semiconductor trade network, and where does Norway sit within it?}

\bigskip
\noindent\textbf{RQ1:} \textit{Which countries occupy structurally central positions in the global semiconductor trade network, and where do potential chokepoints emerge across the front-end and back-end layers?}


\bigskip
\noindent\textbf{RQ2:} \textit{What structural position does Norway occupy across the front-end and back-end segments of the semiconductor value chain, and what does this reveal about its strategic dependencies and leverage?}

\bigskip
\noindent\textbf{RQ3:} \textit{To what extent does political alignment structure the semiconductor trade networks, as measured by UN voting similarity?}

% ─────────────────────────────────────────────────────────────────────────────
\section{Structure of the Thesis}

The thesis proceeds as follows. Chapter~2 establishes the theoretical
framework, connecting the Global Value Chain literature, weaponised
interdependence, and friendshoring. Chapter~3 describes the data sources and
operationalisation. Chapter~4 outlines the two-stage SNA methodology.
Chapter~5 presents the results --- descriptive SNA, centrality analysis,
community detection, and ERGMs. Chapter~6 discusses the findings in relation
to Norway's strategic position. Chapter~7 concludes and notes limitations.
